\chapter{Convergence and the Sobolev Energy}
\label{chap:convergence-sobolev-energy}

The last chapter left a finite-element solution on a partition and a mixed coupling term waiting to be
named. Before the next chapter can lift that term out and call it the single invariant, two things have
to be in place: an account of what happens as the partition is refined toward the continuum, and a way
of measuring the size of a solution that is honest enough to tell genuine structure from nothing. This
chapter supplies both. It shows that the residual of the finite-element solution is squeezed to zero as
the mesh is refined, so the finite answers approach a limit that lives in the completion; and it builds
the energy --- a symmetric, positive-definite form, the Sobolev activity norm --- whose defining virtue
is that it detects zero. The energy will turn out to be the second variation itself, and its
zero-detecting definiteness is what makes the second variation a measuring instrument rather than a
mere quantity.

\section{The squeezed residual}
\label{se:the-squeezed-residual}

Refinement is the lever. A finite-element solution on a coarse partition leaves a residual --- the
amount by which it fails to satisfy the equation exactly --- and refining the partition shrinks that
residual. The construction proves this squeezing directly: as the mesh spacing is driven toward zero,
the pointwise residual is driven toward zero with it, bounded above by a quantity that vanishes in the
limit. The finite answers do not merely get closer to each other; they converge, and the limit they
converge to lives in the completion the construction has carried since its second chapter --- the
Weierstrass completion, the closure in which limits of finite approximations are guaranteed a home. The
contravariant engine of that early chapter is doing the work again here. Each refinement is a finer cut
against the domain, and the sequence of finite solutions is funged toward a single limit point under
the characteristic of agreeing in the limit; the residual squeezed to zero is the cut deciding, against
the whole space of approximations, which one they are all approaching.

A concrete refinement makes the squeezing visible. Solve the problem on a partition of four pieces and
the finite-element solution will miss the true one by some residual --- small, but nonzero, because four
pieces cannot bend to every shape the true solution takes. Halve every piece, to eight, and the residual
roughly quarters: the finer mesh has twice the freedom in each direction and bends closer. Halve again,
to sixteen, and it quarters again. The residual does not wander or stall; it is bounded above by a
quantity that falls with the mesh spacing, and a quantity bounded above by something falling to zero
falls to zero itself. The sequence of finite solutions is therefore Cauchy in the construction's exact
sense --- the crowd of terms agreeing to stop wandering apart --- and the completion is precisely the
home the second chapter built for such a crowd's limit. Nothing here completes the continuum; the
construction only certifies that its finite answers are converging, and names the place they converge
to.

The experiment on refinement makes the mechanism concrete from the other side. What one observer
records as an excess of uncorrelated events over another is a mismatch in refinement depth --- the two
ledgers cannot be aligned without reindexing, and the mismatch must be expressed as a relative ordering
of anchors. Read against the finite-element method, the residual is exactly such a mismatch: it is the
refinement-depth gap between the finite solution and the true one, the events the coarse mesh has not
yet resolved. Driving the residual to zero by refining is closing that gap, aligning the finite
ledger with the continuum's, one refinement at a time. Convergence, in this reading, is not a finite
solution magically becoming exact; it is the disciplined elimination of a refinement-depth mismatch,
and the rate at which it can be eliminated is governed, as ever, by how densely the partition is
allowed to commit its distinctions.

One thing this convergence does not give, and the construction is careful not to claim, is the limit
itself in closed form. What it gives is better: a sequence of finite answers and a guaranteed rate at
which they improve, which is enough to bound how far any one of them sits from the limit. Knowing the
residual quarters with each halving of the mesh, the construction can say not merely that the answers
converge but how fast, and from the rate it can estimate the limit more sharply than any single finite
answer reaches --- the same trick a careful experimentalist uses when she extrapolates a trend rather
than trusting the last data point. The limit is approached on a schedule, and a known schedule is almost
as good as arrival: it tells the construction exactly how much partition it must spend to reach any
accuracy it is asked for.

\section{The energy that detects zero}
\label{se:the-energy-that-detects-zero}

To do anything with these solutions the construction needs to measure them, and the measure it builds
is the energy. Formally it is a symmetric bilinear form --- the Sobolev activity norm --- that takes a
solution and returns a number, its energy-squared, the accumulated activity of its variation. Two
properties make it the right measure. It is symmetric: the energy pairs a solution with itself in a
way that does not depend on which copy is named first. And, the property the whole chapter is built
toward, it detects zero: the energy-squared of a solution is zero if and only if the solution is the
zero solution. Nothing with genuine structure has zero energy, and only genuine emptiness reads as
empty.

It is worth being concrete about what the energy measures. For a solution that varies --- that bends,
that has slope from place to place --- the energy accumulates the square of that variation across the
whole domain, a running total of how much the solution is doing. A flat solution that never changes
contributes nothing to this total; a solution that swings hard contributes a great deal. Squaring is
what makes the contributions add rather than cancel: a downward slope and an upward slope are equally
energetic, because each is squared to a positive number before being summed, so the energy cannot be
talked down to zero by a solution wiggling one way as much as the other. That squaring is also why the
form is symmetric and built as a bilinear pairing: the energy of a solution is the solution paired with
itself, and the same pairing applied to two different solutions measures how much of one lies along the
other. The whole apparatus of angles and lengths --- of one solution being orthogonal to another, of a
solution's size --- comes from this single symmetric pairing, which is why the energy is not just a
number attached to a solution but the geometry the solutions live in.

This is more than a technical nicety, because it is exactly what a measuring instrument must do. An
instrument that reported zero for things that were not there would be worse than useless; the value of
the energy is that its vanishing is reliable testimony of absence. The construction reads the energy as
a genuine norm precisely because it detects zero --- a norm-squared that vanished on something nonzero
would not be a norm at all, only a degenerate form that confuses structure with emptiness. The
experiment on the conservation of energy supplies the physical anchor. For a real field, the energy
assembled from the field's own variation --- its stress and its activity, in the language of the
stress-energy tensor --- is the conserved quantity, the invariant that the dynamics preserve. The
construction's energy is the finite shadow of that: a positive, symmetric measure of a solution's
activity, conserved in the sense that it is the invariant content the second variation will turn out to
be. The energy detecting zero is the construction's guarantee that this invariant is honest --- that
when it says nothing is happening, nothing is.

There is a reason this energy will matter more than an ordinary norm, and it is worth flagging before
the next chapter draws it out. The activity the energy accumulates is not the size of the solution
itself but the size of its \emph{variation} --- how the solution changes, squared and summed. That is
precisely the quantity the eleventh chapter found in the mixed coupling term: the product of two
variations, the second-order response, the thing left over once the first variation has gone flat. The
energy and the second variation are, the construction will show, the same object seen from two sides ---
one built as a norm that detects zero, the other as the curvature of an action at a stationary path.
That they coincide is the hinge the whole part turns on, and the zero-detecting definiteness earned in
this chapter is what equips the second variation to serve as a measure.

\section{Definiteness is earned}
\label{se:definiteness-is-earned}

The zero-detecting property is not free, and the construction is scrupulous about paying for it. A
bilinear form built from variations alone is not automatically positive-definite. It has a nullspace:
the constant and affine modes, solutions whose activity genuinely is zero even though the solutions
themselves are not. A constant function does not vary, so its activity-energy is zero, and a form that
cannot tell a nonzero constant from the zero solution fails the very zero-detection the previous section
demanded. Left to itself the energy is only positive \emph{semi}-definite, blind along its nullspace.

Definiteness has to be earned, and it is earned at the boundary. The construction commits a loaded
boundary anchor --- a boundary condition that pins the solution's value where the domain ends --- and
that single commitment kills the nullspace. With the boundary loaded, a constant mode can no longer
hide: forced to a definite value at the boundary, the only constant consistent with the anchor is the
zero constant, so the constant mode is driven to zero and the affine modes with it. The construction
tanges the offending nullspace by exactly the trait that exposes it --- being unpinned at the boundary
--- and the boundary load removes it. What remains is a genuinely positive-definite form, an energy
that detects zero everywhere because the one place it was blind has been anchored shut. Positive-
definiteness, in this construction, is not an assumption written at the top of the page; it is a
property bought with a boundary condition, and the construction shows the receipt.

The mechanism is worth seeing once in full. Suppose the solution is allowed to be any constant --- a
flat line at any height. Its activity is zero at every point, so its energy is zero, and yet it is not
the zero solution unless its height happens to be zero. An energy that cannot distinguish a flat line at
height seven from the flat line at height zero is blind exactly there, and that blindness is the
nullspace. Now impose the boundary anchor: require that at the edge of the domain the solution take a
fixed value, say zero. A flat line at height seven now violates the anchor --- it reads seven at the
boundary, not zero --- and is no longer admissible. The only constant that survives the anchor is the
one already at zero, which is the zero solution. The nullspace has been emptied of everything but zero,
and the energy, blind nowhere now, detects zero everywhere. One boundary condition, honestly committed,
converts a degenerate form into a genuine norm; the construction never asserts definiteness, it
purchases it.

It is worth naming what the boundary anchor is, physically, because the construction leans on it again
later. A loaded boundary is a place where the domain meets something outside it and exchanges with it
--- a wall held at a fixed temperature, an end pinned to a support, a field clamped at the edge of its
region. Loading the boundary is committing a fact about that exchange, and it is exactly such a
committed boundary fact, in the eighth chapter's sense, that removes the ambiguity a free boundary
leaves. The constant mode is the freedom to add the same nothing everywhere; the boundary load is the
one recorded interaction that says where everywhere is measured from. That a positive-definite energy
requires a committed boundary is not a technical inconvenience but the construction restating, in the
language of forms and nullspaces, a principle it has held since the start: a quantity is definite only
relative to an anchor, and an unanchored measurement measures nothing in particular.

\section{What is demonstrated, and what is assumed}
\label{se:demonstrated-assumed-ch12}

The separation here is clean, and it remains, as in the last two chapters, free of any new assumption.

What is demonstrated is convergence and a measure. The finite-element residual is proved to be squeezed
to zero as the mesh refines, with the limit residing in the completion. The energy is built as a
symmetric bilinear form and proved to detect zero --- its energy-squared vanishes exactly on the zero
solution --- once definiteness is earned. And definiteness is earned, not assumed: the bilinear form's
constant-and-affine nullspace is proved to be killed by a loaded boundary anchor, which forces the
constant mode to zero. Each of these is a theorem about finite objects, and none of them reaches for
the oracle of the tenth chapter; convergence and the energy stand on construction alone.

It is worth pausing on how little this chapter cost, because the contrast is the point. The whole
apparatus --- a proof of convergence, a symmetric energy, a theorem that it detects zero, and a second
theorem that a boundary load earns the definiteness --- is built and checked without a single new
assumption. The tenth chapter's oracle is not invoked here; convergence is a squeezing the construction
performs, and definiteness is a purchase it makes at the boundary. After the keystone, the book has
returned to solid ground, and it stays there: everything from here to the second variation's naming is
construction on top of the one assumption already declared, and the energy this chapter earns is the
instrument the naming will use.

\begin{quote}
\emph{A note on the labels.} Two bridges carry physical weight here. Calling the symmetric form an
\emph{energy} borrows the word from physics for what is, in itself, an accumulated measure of
variation; the conservation-of-energy reading is earned only to the degree the finite measure tracks
the physical invariant, which the next chapter will make precise. And the \emph{completion} in which the
refined solutions find their limit is the same modeling bridge met at the number tower: the continuum is
approached, never constructed, and to say the limit ``lives in the completion'' is to assert that the
approach has a destination, an assertion licensed by the squeezing but not identical to having reached
it.
\end{quote}

What is assumed is therefore only those bridges and the standing oracle of stationarity --- nothing
new. With a notion of convergence in hand and an energy that reliably detects zero, the construction has
exactly the equipment the next chapter needs. It can take the mixed coupling term that the splines left
waiting, recognize it as this very energy, and prove the result the whole part has been named for: that
the second variation is the single invariant, the one quantity whose vanishing certifies a stationary
path and whose magnitude measures everything a measurement can detect about how that path responds.
